\section{Contrôler le repo à partir des contrôles forward et dividendes}\label{sec-controle}

\begin{objectif}
(O2) Répondre à la question \og j'ai le contrôle forward, j'ai le contrôle dividendes : combien vaut le contrôle
repo ? \fg{}, et mesurer, portefeuille par portefeuille, ce que la répartition de Totem entre repo et dividendes change
au contrôle.
\end{objectif}

\subsection{Lignes de contrôle}
On regroupe les taux de repo aux piliers et les niveaux de dividendes par période dans
$\theta=(q, \ell)\in\R^N\times\R^m$. Le spot, la courbe des taux et le profil d'allocation sont fixés, de sorte que les
forwards aux piliers forment une fonction régulière $F(\theta)\in\R^N$. D'après la \cref{prop-bootstrap}, pour des
forwards $F$ et des niveaux $\ell$ donnés, il existe un unique repo $q=\psi(F, \ell)$ tel que
$F\bigl(\psi(F, \ell), \ell\bigr)=F$ : c'est l'\emph{application repo implicite}.

Les marquages FO sont $\theta^{\fo}=(q^{\fo}, \ell^{\fo})$. Totem fournit des forwards $F^{\tot}$ et des niveaux
$\ell^{\tot}$, d'où $q^{\tot}=\psi(F^{\tot}, \ell^{\tot})$ et $\hat q=\psi(F^{\tot}, \ell^{\fo})$. Un portefeuille a une
valeur $V(\theta)$ et des sensibilités FO $s^q=\partial_qV(\theta^{\fo})\in\R^N$ et
$s^\ell=\partial_\ell V(\theta^{\fo})\in\R^m$. On note enfin
\[
\Delta\ell=\ell^{\tot}-\ell^{\fo}, \qquad\Delta q=q^{\tot}-q^{\fo}, \qquad\Delta\hat q=\hat q-q^{\fo}.
\]

\begin{definition}[Lignes de contrôle]\label{def-lines}
\[
\text{ligne dividendes}=\sprod{s^\ell}{\Delta\ell}, \qquad\text{ligne repo}=\sprod{s^q}{\Delta q}, \qquad
\text{ligne forward}=\sprod{s^q}{\Delta\hat q}.
\]
Le \emph{total} est la somme des lignes dividendes et repo. La ligne forward est la ligne repo que l'on obtiendrait en
gardant les dividendes du FO : elle ne dépend de Totem qu'à travers ses forwards.
\end{definition}

\begin{definition}[Portefeuille linéaire en forward]\label{def-fwdlinear}
Un portefeuille est \emph{linéaire en forward} si $V(\theta)=\varphi\bigl(F(\theta)\bigr)$ pour une fonction $\varphi$
de classe $C^2$ : sa valeur ne dépend du repo et des dividendes qu'à travers les forwards. C'est le cas des options
européennes valorisées avec une nappe de volatilité qui ne dépend pas de la répartition, des forwards et des futures.
\end{definition}

\subsection{La jacobienne du repo implicite}
\begin{lemma}[Jacobienne]\label{lem-jacobian}
La matrice $J=\partial_\ell\psi(F, \ell)\in\M_{N,m}(\R)$ vaut $J=-(\partial_qF)^{-1}\partial_\ell F$. La matrice
$\partial_qF$ est triangulaire inférieure, de termes diagonaux $\partial F_k/\partial q_k\le-T_kF_k<0$, donc
inversible. Au premier ordre,
\begin{equation}\label{eq-jacobian}
J_{kp}\approx-\frac{a_{kp}}{T_k\,(S_0-D_k)} .
\end{equation}
\end{lemma}
\begin{proof}
On dérive $F\bigl(\psi(F, \ell), \ell\bigr)=F$ par rapport à $\ell$ : $\partial_qF\,J+\partial_\ell F=0$. Avec un repo
forward constant entre piliers, $F_k$ ne dépend que de $q_1, \dots, q_k$, d'où la forme triangulaire ; le terme
diagonal vaut $T_k\,\mathcal F_k'(q_kT_k)\le-T_kF_k$ d'après la \cref{prop-bootstrap}. Pour le premier ordre, on fixe
les forwards dans \eqref{eq-additivity} : $T_k\,\delta q_k=-\delta D_k/(S_0-D_k)=-a_{kp}\,\delta\ell_p/(S_0-D_k)$.
\end{proof}

\subsection{L'identité de contrôle}
\begin{theorem}[Identité de contrôle]\label{thm-control}
Pour tout portefeuille,
\begin{equation}\label{eq-control}
\text{ligne dividendes}+\text{ligne repo}=\text{ligne forward}+\sprod{\eta}{\Delta\ell}+O\bigl(\norm{\Delta}^2\bigr),
\end{equation}
où $\eta=s^\ell+\tr{J}\,s^q\in\R^m$, $J$ est évaluée en $(F^{\tot}, \ell^{\fo})$ et $\norm{\Delta}=\norm{\Delta\ell}+\norm{\Delta\hat q}$. On appelle
$\sprod{\eta}{\Delta\ell}$ la \emph{ligne non-forward} et $\eta$ l'\emph{exposition non-forward}. Si le portefeuille
est linéaire en forward, $\eta=O(\norm{\Delta\hat q})$ : le total est égal à la ligne forward au second ordre près, et
ne dépend donc pas des dividendes Totem au premier ordre.
\end{theorem}
\begin{proof}
Un développement de $\psi$ en $\ell$ donne $\psi(F^{\tot}, \ell^{\tot})=\psi(F^{\tot}, \ell^{\fo})+J\Delta\ell
+O(\norm{\Delta\ell}^2)$, soit $\Delta q=\Delta\hat q+J\Delta\ell+O(\norm{\Delta\ell}^2)$. Donc
\[
\sprod{s^\ell}{\Delta\ell}+\sprod{s^q}{\Delta q}=\sprod{s^q}{\Delta\hat q}+\sprod{s^\ell+\tr{J}s^q}{\Delta\ell}
+O(\norm{\Delta}^2).
\]
Si $V=\varphi\circ F$, alors $\tr{s^q}=\nabla\varphi\,\partial_qF(\theta^{\fo})$ et
$\tr{s^\ell}=\nabla\varphi\,\partial_\ell F(\theta^{\fo})$, d'où
$\tr{\eta}=\nabla\varphi\,\bigl[\partial_\ell F+\partial_qF\,J\bigr](\theta^{\fo})$. Au point
$\theta'=(\hat q, \ell^{\fo})$, le crochet est nul d'après le \cref{lem-jacobian} ; comme
$\theta'-\theta^{\fo}=(\Delta\hat q, 0)$ et que les dérivées sont lipschitziennes, le crochet en $\theta^{\fo}$ est un
$O(\norm{\Delta\hat q})$, et $\sprod{\eta}{\Delta\ell}=O(\norm{\Delta}^2)$.
\end{proof}

\begin{reponse}
En réécrivant \eqref{eq-control} :
\[
\boxed{\;\text{ligne repo}=\text{ligne forward}-\text{ligne dividendes}+\text{ligne non-forward}\;}
\]
au second ordre près.
Le contrôle repo se déduit du contrôle forward et du contrôle dividendes. Pour un portefeuille linéaire en forward, la
ligne non-forward est nulle : le contrôle repo est exactement la différence des deux autres, et le total (dividendes
$+$ repo) ne dépend pas des contributeurs dividendes de Totem. Le panel dividendes n'intervient qu'à travers
l'exposition non-forward $\eta=s^\ell+\tr{J}s^q$, qui se calcule à partir des seules sensibilités FO et de la
jacobienne du bootstrap, sans revalorisation.
\end{reponse}

Le résultat est même exact en réévaluation complète.

\begin{theorem}[Exactitude en réévaluation complète]\label{thm-exact}
Si le portefeuille est linéaire en forward et que toutes ses maturités sont des piliers Totem, le contrôle en
réévaluation complète $V(\theta^{\tot})-V(\theta^{\fo})=\varphi(F^{\tot})-\varphi\bigl(F(\theta^{\fo})\bigr)$, avec
$\theta^{\tot}=(q^{\tot}, \ell^{\tot})$, ne dépend pas des dividendes Totem $\ell^{\tot}$.
\end{theorem}
\begin{proof}
Par construction du repo implicite, $F(\theta^{\tot})=F\bigl(\psi(F^{\tot}, \ell^{\tot}), \ell^{\tot}\bigr)=F^{\tot}$, quel
que soit $\ell^{\tot}$.
\end{proof}

\begin{remark}[Maturités entre deux piliers]
Pour une maturité située entre deux piliers, le forward dépend du portage interpolé et des dividendes détachés entre
les deux : une dépendance résiduelle à $\ell^{\tot}$ subsiste. Elle est de l'ordre des dividendes détachés entre la
maturité et les piliers voisins, c'est-à-dire exactement les effets de date et de profil de la
\cref{prop-signatures}, que surveille le détecteur. Le terme du second ordre de \eqref{eq-control} est explicité en
\cref{app-second}.
\end{remark}

\subsection{Exemples}
\begin{example}[Un call]\label{ex-call}
Un call à la monnaie à un an : spot $100$, taux \pct{3}, volatilité \pct{25}, dividende FO $D^{\fo}=\num{3.5}$, repo FO
$60$\,pb. Le forward Totem est inférieur de $20$\,pb à celui du FO, et l'on essaie trois panels dividendes :
\begin{center}\input{generated/\lang/tab_cancel}\end{center}
Le repo implicite varie de $125$\,pb d'un panel à l'autre et chaque ligne d'environ $\num{0.3}$, mais le total ne
bouge pas, et il est égal à la réévaluation complète.
\end{example}

\begin{example}[Trois portefeuilles]\label{ex-books}
Le pricer du prototype valorise trois portefeuilles sur onze piliers de repo et six périodes de dividendes :
A (calls, puts et un forward), B (A plus un swap de dividendes 2028 vendu) et C (A plus une commission de prêt de titres
à un an). Les forwards Totem sont inférieurs de $20$\,pb à ceux du FO, et le panel dividendes Totem décale les
périodes 2027 et 2028 de $\Delta\ell\in\{\num{-0.4}\,;\,0\,;\,\num{+0.5}\}$ :
\begin{center}\input{generated/\lang/tab_books}\end{center}
Pour A, le total reste à \vBookATotal{} (la ligne forward) alors que chaque ligne varie de plus de six unités. Pour B et
C, le total bouge, exactement de la ligne non-forward : le swap de dividendes est une vraie exposition aux dividendes,
la commission de prêt une vraie exposition au repo. La réévaluation complète confirme chaque ligne
(\cref{fig-cancellation}).
\end{example}

\begin{figure}[ht]
\centering\includegraphics[width=\linewidth]{figures/\lang/cancellation.pdf}
\caption{Lignes dividendes et repo, et leur total, quand le panel dividendes Totem se déplace (portefeuille A à
gauche, B à droite). Points : réévaluation complète. Pour un portefeuille linéaire en forward, le total est plat
(\cref{thm-control,thm-exact}) ; le swap de dividendes ajoute une vraie exposition aux dividendes, dont la ligne est
exactement $\sprod{\eta}{\Delta\ell}$.}
\label{fig-cancellation}
\end{figure}

\subsection{Ce que la répartition coûte à un portefeuille}
\begin{proposition}[Matérialité de la répartition]\label{prop-materiality}
On suppose que l'erreur de consensus sur les niveaux de dividendes $\Delta\ell$ a une matrice de covariance
$\Gamma^\ell$, les forwards Totem étant fixés. Au premier ordre, la variance du total vaut $\tr{\eta}\,\Gamma^\ell\eta$,
contre $\tr{(s^\ell)}\,\Gamma^\ell s^\ell$ pour la ligne dividendes seule. La part de l'incertitude de la ligne dividendes
compensée par la ligne repo vaut donc
\[
1-\sqrt{\frac{\tr{\eta}\,\Gamma^\ell\eta}{\tr{(s^\ell)}\,\Gamma^\ell s^\ell}} .
\]
\end{proposition}
\begin{proof}
D'après le \cref{thm-control}, le total est la ligne forward, qui ne dépend pas de $\ell^{\tot}$, plus
$\sprod{\eta}{\Delta\ell}$.
\end{proof}

\begin{example}[Matérialité]\label{ex-materiality}
Avec une dispersion de consensus de \pct{5} de chaque période, dix contributeurs et une corrélation $e^{-|c_p-c_{p'}|/2}$
entre périodes de milieux $c_p$ et $c_{p'}$ :
\begin{center}\input{generated/\lang/tab_materiality}\end{center}
\vMatA{} de l'incertitude de la ligne dividendes se compense pour le portefeuille A, \vMatB{} pour B et \vMatC{} pour C.
\end{example}

\subsection{Quand la répartition compte}
L'exposition non-forward $\eta$ n'est pas nulle pour les swaps et futures de dividendes (exposition pure aux
dividendes), pour les commissions de prêt-emprunt de titres et les jambes de financement des swaps de rendement total
(exposition pure au repo), pour l'exercice américain (la frontière d'exercice anticipé dépend séparément des
dividendes et du repo) et pour les exotiques valorisés avec une dynamique de dividendes discrets. Pour ces
portefeuilles, la répartition doit être contrôlée directement, et la qualité de la courbe de repo implicite compte :
c'est l'objet de la \cref{sec-detection}.

\begin{retenir}
\begin{itemize}[leftmargin=12pt]
\item Contrôler la \emph{ligne forward} (forwards Totem, dividendes FO) et la \emph{ligne non-forward}
$\sprod{\eta}{\Delta\ell}$ ; la seconde n'a besoin des dividendes Totem que là où $\eta\ne0$.
\item Des lignes dividendes et repo importantes et de signes opposés, avec un total faible, sont un artefact de
répartition, pas une erreur de marquage.
\item En valorisation prudente~\cite{eu2016}, agréger les incertitudes par paramètre sans compensation compte deux
fois un artefact de répartition ; prendre le forward comme paramètre là où $\eta\approx0$ supprime ce double compte.
\end{itemize}
\end{retenir}
